\chapter{Continued Fractions and Diophantine Approximation}\label{continued-fractions}
\chapterstatus{planned}

How well real numbers can be approximated by rationals, and what continued fractions say about it.

\noindent\textbf{Prerequisites.} Chapter~\ref{number-theory} (Elementary Number Theory) and Chapter~\ref{real-analysis} (Real Analysis).

\noindent\textbf{Planned contents.}
\begin{enumerate}
\item Finite and infinite continued fractions, convergents and their recurrences; the Euclidean algorithm again.
\item Best approximations; Dirichlet's approximation theorem, Hurwitz's theorem and the constant $\sqrt{5}$; Farey sequences and Ford circles.
\item Quadratic irrationals: Euler's and Lagrange's theorems on periodic continued fractions; reduced quadratic irrationals (Galois).
\item Pell's equation $x^2 - dy^2 = \pm 1$: existence of solutions, the fundamental solution, the group of solutions.
\item Liouville's theorem and Liouville numbers; transcendence of $e$ and of $\pi$ (Hermite, Lindemann); the Thue--Siegel--Roth theorem (statement).
\item Equidistribution modulo $1$: Weyl's criterion and $n\alpha$; the three-gap theorem.
\item Exercises with solutions.
\end{enumerate}
